% 第 4 章 伽马射线谱学离线数据处理与物理量提取方法

\chapter{伽马射线谱学离线数据处理与物理量提取方法}
\label{ch:data_analysis_methods}

\section{多重符合矩阵构建与开门本底扣除技术}
\label{sec:coincidence_matrix}

在重离子熔合蒸发反应中，出射道通常同时开放数个中子蒸发道、带电粒子道以及巨裂变副反应，靶室中瞬间产生数十种不同的剩余同位素核。如果仅仅测量未经符合筛选的单能谱（Singles Spectrum），成千上万条来自不同核素的 $\gamma$ 射线在能谱上严重重叠挤压，根本无法辨析任何一条能带的级联序列。

现代在束核谱学的基石是\term{多重符合技术}（Coincidence Spectroscopy）。其物理基础是：处于同一级联衰变能带中的多个激发能级，其退激寿命通常在皮秒到纳秒量级；当原子核从某一高自旋态逐级向下退激时，发射出的数个 $\gamma$ 光子在微观时间尺度上几乎是同时发射的。

\subsection{二维符合矩阵与时间符合窗}

通过对实验记录的 List-Mode 离散事件流进行扫描，在设定的纳秒级时间符合窗（通常 $\Delta t \le 30 \sim 50\text{ ns}$）内，如果探测阵列中有两个或两个以上的独立探测单元同时响应，系统将其判定为一次有效符合事件。

对于双重符合事件，离线程序构建一个 $4096 \times 4096$ 或 $8192 \times 8192$ 通道的对称二维矩阵 $M(E_1, E_2)$。矩阵元 $M(i, j)$ 的数值代表第一个探测器记录到能量通道 $i$、同时第二个探测器记录到能量通道 $j$ 的总符合计数。为了完全消除探测器寻址引起的非对称性，离线处理中普遍对矩阵进行对称化操作：$M_{\text{sym}}(i, j) = M(i, j) + M(j, i)$。

\subsection{开门技术（Gating）与二维全局本底扣除}

所谓\term{开门}（Gating），就是在二维矩阵的一个轴上指定某一条确定属于目标能带的 $\gamma$ 跃迁能量峰作为“门道窗口”（$E_{\text{gate}} \pm \Delta E$），然后沿另一个能量轴切出一条一维谱线。这样切出的能谱中，理论上只保留了与这扇“门”在物理上存在前后级联符合关系的 $\gamma$ 跃迁。

然而，门道窗口内不仅包含全能峰计数，还必然包含了下方深沉的康普顿散射本底以及其它偶然杂质峰。为了获得真实的物理级联强度，必须进行严格的\term{本底扣除}（Background Subtraction）。

在国际通用的 Radware 分析程序（D. C. Radford 发展）与 P. J. Papanicolas 算法中，采用二维全局本底分解模型：
\begin{equation}
    B(E_1, E_2) = \frac{T(E_1) \cdot B(E_2) + B(E_1) \cdot T(E_2) - B(E_1) \cdot B(E_2)}{T_{\text{total}}}
    \label{eq:matrix_background}
\end{equation}
其中 $T(E)$ 为单轴总投影谱，$B(E)$ 为平滑拟合的一维连续本底谱。

通过从门峰切谱中扣除归一化的门边本底谱（Sideband Gates），能够极其干净地消除康普顿连续散射与随机偶发符合事件，使信噪比（S/N）提高数十倍，从而清晰辨识强度仅占最强主线千分之几的高自旋弱跃迁。

\section{方向角分布（ADO）与方向角关联（DCO）理论体系}
\label{sec:dco_formalism}

建立了各 $\gamma$ 射线之间的符合关系后，最核心的物理任务是确定每一个核能级的\term{自旋与宇称}（$I^\pi$），以及各跃迁的\term{多极性质}（电四极 $E2$、磁偶极 $M1$ 或电偶极 $E1$）。

\subsection{核取向态与方向角分布展开公式}

在重离子熔合反应中，由于入射束流具有沿束流轴（$z$ 轴）的高度定向动量，带入的角动量 $\vec{\ell}$ 严格垂直于束流方向。形成的复合核其角动量高度对齐在垂直于束流的平面内，磁子态分布呈现高度取向（$m \approx 0$ 高度集中，即核完全对齐，Nuclear Alignment）。

根据核辐射场多极展开理论，定向核发射光子的角分布不是各向同性的，而是严格依赖于探测器相对束流轴的夹角 $\theta$。方向角分布（Angular Distribution of Oriented nuclei, ADO）的一般代数形式为：
\begin{equation}
    W(\theta) = A_0 \left[ 1 + a_2 P_2(\cos\theta) + a_4 P_4(\cos\theta) \right]
    \label{eq:angular_distribution}
\end{equation}
其中 $P_2(x) = \frac{1}{2}(3x^2 - 1)$ 与 $P_4(x) = \frac{1}{8}(35x^4 - 30x^2 + 3)$ 为 Legendre 多项式。

展开系数 $a_k$ 在微观上由两部分相乘决定：
\begin{equation}
    a_k = B_k(I_i) \cdot F_k(I_f, L_1, L_2, I_i)
\end{equation}
其中：
\begin{itemize}
    \item $B_k(I_i)$：表征初态核取向程度的统计取向张量（Orientation Tensor），通常假定子态呈高斯分布，取向宽度参数为 $\sigma/I$；
    \item $F_k$：完全由角动量耦合 Racah 代数与 Wigner $6j$ 符号决定的纯几何系数（Yamazaki $F$-系数）。
\end{itemize}

统计取向张量由磁子态密度矩阵对角元 $\rho_{mm}(I)$ 精确定义：
\begin{equation}
    B_k(I) = \sqrt{2I+1} \sum_{m=-I}^I (-1)^{I-m} C_{I m I -m}^{k 0} \, P(m)
    \label{eq:orientation_tensor}
\end{equation}
在重离子在束反应中，磁子态占据几率普遍遵从以 $m=0$ 为中心的高斯分布模型：
\begin{equation}
    P(m) = \frac{\exp\left( -\frac{m^2}{2\sigma^2} \right)}{\sum_{m'=-I}^I \exp\left( -\frac{m'^2}{2\sigma^2} \right)}
    \label{eq:gaussian_substates}
\end{equation}
其中无量纲参量 $\sigma/I$ 表征取向宽度。在典型熔合反应中，$\sigma/I \approx 0.25 \sim 0.35$。当核完全对齐（完全局域在 $m=0$）时，取向因子达到理论极大；随后未被观测到的高能级联过渡使取向发生轻度去取向（Deorientation），在角分布拟合中通常引入衰减因子 $\alpha_k$ 进行经验修正。

Yamazaki $F_k$ 系数的精确代数封闭解析式为：
\begin{equation}
    F_k(I_f, L_1, L_2, I_i) = (-1)^{I_f - I_i - 1} \sqrt{(2L_1+1)(2L_2+1)(2I_i+1)} \, C_{L_1 1 L_2 -1}^{k 0} \, W(I_i I_i L_1 L_2; k I_f)
    \label{eq:yamazaki_formula}
\end{equation}
其中 $C_{L_1 1 L_2 -1}^{k 0}$ 为 Clebsch-Gordan 系数，$W(I_i I_i L_1 L_2; k I_f)$ 为 Racah 算符系数。该式完全摆脱了对经验查表的依赖，实现了在任意自旋与任意多极性混合下的严格计算机数值求解。

\subsection{方向角关联（DCO）比值与典型判据}

在多探测器阵列中，探测器分布在不同的极角 $\theta$ 上。为了规避单谱拟合对束流电荷积分与探测器死时间的严苛依赖，离线处理普遍采用\term{方向角关联比值}（Directional Correlation from Oriented nuclei, DCO Ratio）。

对于两级联跃迁 $\gamma_1$（待测跃迁）与 $\gamma_2$（门道参考跃迁），DCO 比值定义为：
\begin{equation}
    R_{\text{DCO}} = \frac{I_{\gamma_1}(\theta_1; \text{gate at } \theta_2)}{I_{\gamma_1}(\theta_2; \text{gate at } \theta_1)} \cdot \frac{\epsilon(\theta_2, E_1)}{\epsilon(\theta_1, E_1)}
    \label{eq:dco_ratio}
\end{equation}
在实际阵列配置中，通常选择一组位于前后小角度的探测器（如 $\theta_1 \approx 30^\circ$ 或 $150^\circ$）与一组位于中间垂直角度的探测器（如 $\theta_2 \approx 90^\circ$）构成探测器对。

如果门道参考跃迁 $\gamma_2$ 是一条确定的已知\term{拉伸电四极跃迁}（Stretched $E2$，即 $I \to I-2$）：
\begin{itemize}
    \item 若待测跃迁 $\gamma_1$ 是纯拉伸四极跃迁（Stretched $E2$）：由于前后向计数高于垂直向计数，$R_{\text{DCO}} \approx 1.0$；
    \item 若待测跃迁 $\gamma_1$ 是纯拉伸偶极跃迁（Stretched Dipole $\Delta I = 1$，如纯 $E1$ 或纯 $M1$）：垂直向计数高于前后向计数，因此 $R_{\text{DCO}} \approx 0.5 \sim 0.6$；
    \item 若待测跃迁 $\gamma_1$ 是非拉伸偶极跃迁（Non-stretched Dipole $\Delta I = 0$）：角分布再次倒转，$R_{\text{DCO}} \approx 1.0 \sim 1.1$。
\end{itemize}

\section{多极混合比 $\delta(E2/M1)$ 严格代数双实根反演}
\label{sec:mixing_ratio_dual_root}

在很多学术论文中，研究者常常教条地套用上述经验范围，只要测得 $R_{\text{DCO}} \approx 0.6$ 就盲目宣称是偶极跃迁，测得 $R_{\text{DCO}} \approx 1.0$ 就宣称是四极跃迁。**这在物理方法上存在极大的漏洞与隐患！**

在低能核结构中，绝大多数 $\Delta I = 1$ 的级联跃迁（特别是在同核带内或旋称伙伴带之间的跃迁）都不是纯跃迁，而是磁偶极（$M1$）与电四极（$E2$）强烈混合的杂化跃迁。

\subsection{混合比 $\delta$ 的数学定义与二次有理分式}

多极混合比 $\delta$（Multipole Mixing Ratio）定义为电四极约化矩阵元与磁偶极约化矩阵元的比值：
\begin{equation}
    \delta(E2/M1) = \frac{\langle I_f || \hat{T}(E2) || I_i \rangle}{\langle I_f || \hat{T}(M1) || I_i \rangle}
    \label{eq:mixing_ratio_def}
\end{equation}
此时，跃迁光子中含有比例为 $\frac{\delta^2}{1 + \delta^2}$ 的四极成分，以及 $\frac{1}{1 + \delta^2}$ 的偶极成分。

角分布系数不仅包含纯项，更包含了偶极算符与四极算符之间的量子干涉交叉项：
\begin{equation}
    a_k(\delta) = \frac{F_k(I_f, 1, 1, I_i) + 2\delta \, F_k(I_f, 1, 2, I_i) + \delta^2 F_k(I_f, 2, 2, I_i)}{1 + \delta^2}
    \label{eq:ak_mixing}
\end{equation}
因此，理论上预测的 $R_{\text{DCO}}$ 是混合比 $\delta$ 的\textbf{二次有理分式}：
\begin{equation}
    R_{\text{DCO}}(\delta) = \frac{\alpha_0 + \alpha_1 \delta + \alpha_2 \delta^2}{\beta_0 + \beta_1 \delta + \beta_2 \delta^2}
    \label{eq:dco_quadratic}
\end{equation}

\subsection{代数双实根退化现象（Dual-Root Degeneracy）}

将实验实测的数值 $R_{\text{exp}}$ 带入方程~(\ref{eq:dco_quadratic})，整理后得到关于 $\delta$ 的标准一元二次方程：
\begin{equation}
    A(R_{\text{exp}}) \delta^2 + B(R_{\text{exp}}) \delta + C(R_{\text{exp}}) = 0
    \label{eq:quadratic_equation}
\end{equation}
其中二次项、一次项与常数项系数分别为：
\begin{equation}
    A = \alpha_2 - R_{\text{exp}} \beta_2, \quad B = \alpha_1 - R_{\text{exp}} \beta_1, \quad C = \alpha_0 - R_{\text{exp}} \beta_0
\end{equation}

根据一元二次方程求根公式，混合比的两个解析解为：
\begin{equation}
    \delta_{1, 2} = \frac{-B \pm \sqrt{B^2 - 4AC}}{2A}
    \label{eq:dual_root_solution}
\end{equation}

对于物理上绝大多数测量值，该方程的判别式严格满足 $\Delta > 0$！这意味着：**单个方向角关联测量值 $R_{\text{DCO}}$ 在数学上必然对应着两个完全合法的实数根解**（图~\ref{fig:dual_roots}）：
\begin{itemize}
    \item 一个大混合比根（通常 $|\delta_1| > 1.5$）：四极跃迁成分占优（$> 70\%$）；
    \item 一个小混合比根（通常 $|\delta_2| < 0.4$）：偶极跃迁成分占优（$> 85\%$）。
\end{itemize}

\begin{figure}[htbp]
    \centering
    \begin{tikzpicture}[scale=0.95]
        \draw[thick, ->] (-4,0) -- (4,0) node[right] {混合比 $\delta$};
        \draw[thick, ->] (0,-0.5) -- (0,3.5) node[above] {$R_{\text{DCO}}$};

        % DCO 曲线
        \draw[line width=1.3pt, blue] (-3.8,2.8) .. controls (-2.5,2.9) and (-1.5,1.2) .. (-0.8,0.6) .. controls (-0.3,0.3) and (0.5,0.4) .. (1.2,0.9) .. controls (2.0,1.8) and (3.0,2.6) .. (3.8,2.8);

        % 实验水平线
        \draw[dashed, red, line width=1.1pt] (-3.8,1.2) -- (3.8,1.2) node[right] {$R_{\text{exp}} \pm \sigma$};

        % 双根点
        \filldraw[black] (-1.35,1.2) circle (2.5pt) node[above left] {解 $\delta_1 \approx -1.5$ (大根)};
        \filldraw[black] (1.45,1.2) circle (2.5pt) node[above right] {解 $\delta_2 \approx +0.3$ (小根)};
        \draw[dotted, gray, thick] (-1.35,1.2) -- (-1.35,0) node[below] {$\delta_1$};
        \draw[dotted, gray, thick] (1.45,1.2) -- (1.45,0) node[below] {$\delta_2$};
    \end{tikzpicture}
    \caption{方向角关联方程的一元二次代数结构诱发多极混合比双实根退化示意图}
    \label{fig:dual_roots}
\end{figure}

在所谓“横向摆动运动（TW）”的国际争议中，很多论文之所以能够宣称发现了摆动带，其关键原因就是作者在拟合双根时，**主观选择并仅采纳了大混合比根 $\delta_1$**，从而计算出巨大的 $B(E2)_{\text{out}}$。如果采纳同样自洽的小实根 $\delta_2$，跃迁本质上只是普通的磁偶极退激，摆动假设立即彻底崩溃！

因此，严谨的在束谱学研究规范要求：**绝不允许在没有独立实验证据的前提下主观选根！**必须同时汇报两个实数根，并借助下一节介绍的康普顿散射线偏振测量或内转换系数进行决定性仲裁。

\section{康普顿散射线偏振度测定与电/磁宇称确认}
\label{sec:linear_polarization}

方向角分布和 DCO 只能确定跃迁的多极阶数 $L$（如偶极 $L=1$ 或四极 $L=2$），但无法直接区分该跃迁究竟是\textbf{电跃迁}（Electric, $E$）还是\textbf{磁跃迁}（Magnetic, $M$），因为在相同的角动量几何下，同阶的电与磁跃迁角分布形式完全相同，两者的唯一区别在于辐射电场矢量 $\vec{E}$ 与磁场矢量 $\vec{B}$ 的空间偏振方向互换了 $90^\circ$。

要最终确立能级的宇称变化（$\pi_i \pi_f = (-1)^L$ 对应电跃迁；$\pi_i \pi_f = (-1)^{L+1}$ 对应磁跃迁），必须测量发射光子的\term{线偏振度}（Linear Polarization）。

\subsection{Klein-Nishina 截面与康普顿极化仪}

根据量子电动力学 Klein-Nishina 康普顿散射微分截面公式：
\begin{equation}
    \frac{d\sigma}{d\Omega} = \frac{r_0^2}{2} \left( \frac{E'}{E} \right)^2 \left[ \frac{E'}{E} + \frac{E}{E'} - 2 \sin^2\theta_{\text{sc}} \cos^2\phi_{\text{sc}} \right]
    \label{eq:klein_nishina}
\end{equation}
其中 $\theta_{\text{sc}}$ 为散射角，$\phi_{\text{sc}}$ 为散射方向相对入射光子线偏振电矢量 $\vec{E}$ 的方位角。

方程~(\ref{eq:klein_nishina})揭示了一个极其重要的微观物理事实：**康普顿散射光子具有极强的物理倾向，垂直于入射电矢量 $\vec{E}$ 的方向散射（$\phi_{\text{sc}} = 90^\circ$ 处截面最大，$\phi_{\text{sc}} = 0^\circ$ 处截面最小）**。

在现代探测阵列中，紧密拼接的 Clover 探测器（由 4 个独立的同轴晶体拼接在一个制冷机内）或先进跟踪阵列（AGATA/GRETINA），是天然的高灵敏度\term{康普顿散射极化仪}（Compton Polarimeter）。当一个光子在一个晶体中发生散射并被相邻晶体吸收时，两晶体连线与反应平面的几何夹角直接记录了散射方向。

\subsection{实验不对称度与宇称判据}

定义散射光子垂直于束流与探测器构成的反应平面的计数为 $N_\perp$，平行于反应平面的计数为 $N_\parallel$。实验测量的不对称度为：
\begin{equation}
    \Delta = \frac{N_\perp - a N_\parallel}{N_\perp + a N_\parallel}
\end{equation}
其中 $a$ 为极化仪内部晶体本征效率的几何标定校正系数。

绝对线偏振度 $P$ 表达为：
\begin{equation}
    P = \frac{\Delta}{Q(E_\gamma)}
\end{equation}
其中 $Q(E_\gamma)$ 为极化仪在对应能量下的偏振灵敏度因子（Polarization Sensitivity，通常在 $0.15 \sim 0.35$ 之间）。

对于理想无体积点状极化仪，在 $90^\circ$ 康普顿散射角下，极化灵敏度 $Q_0$ 由 Klein-Nishina 截面解析推导为：
\begin{equation}
    Q_0(E_\gamma) = \frac{1 + \frac{E_\gamma}{m_e c^2}}{1 + \frac{E_\gamma}{m_e c^2} + \left(\frac{E_\gamma}{m_e c^2}\right)^2}
    \label{eq:polarimeter_sensitivity}
\end{equation}
在低能区（$E_\gamma \to 0$），$Q_0 \to 1.0$ 达到完美偏振灵敏度；随着能量上升，多重散射几率增大，$Q$ 因子单调下降，实际探测器阵列必须通过已知纯电/纯磁跃迁进行高精度的实验连续能量刻度标定。

结合理论角分布推导，电磁性质判定定则极其明晰：
\begin{itemize}
    \item 对于拉伸偶极跃迁（$\Delta I = 1$）：若测量得到 $P > 0$，确认为无宇称改变的纯\textbf{磁偶极跃迁（$M1$）}；若 $P < 0$，确认为改变宇称的纯\textbf{电偶极跃迁（$E1$）}；
    \item 对于拉伸四极跃迁（$\Delta I = 2$）：若测量得到 $P > 0$，确认为不改变宇称的纯\textbf{电四极跃迁（$E2$）}。
\end{itemize}
偏振测量是彻底砸碎“双实根简并”、最终裁定电偶极还是电四极跃迁的不可逾越的法官。

\section{激发态皮秒能级寿命测量与绝对跃迁几率}
\label{sec:lifetime_measurements}

能级能量、DCO 与偏振测量只能揭示核波函数的空间几何构型与多极对称性，无法给出核四极形变的绝对大小。要获得原子核的绝对内禀电四极矩 $Q_0$ 以及约化四极跃迁几率 $B(E2)$，**必须测定激发态的寿命 $\tau$（或半衰期 $T_{1/2} = \tau \ln 2$）**。

根据电磁跃迁理论，态跃迁几率与绝对跃迁矩阵元严格绑定：
\begin{equation}
    B(E2; I \to I-2) = \frac{0.08156}{\tau \cdot E_\gamma^5 \cdot (1 + \alpha)} \quad [e^2\text{b}^2]
    \label{eq:be2_lifetime}
\end{equation}
高自旋态能级的寿命通常在飞秒（$10^{-15}\text{ s}$）到皮秒（$10^{-12}\text{ s}$）量级，远快于任何常规电子学的响应时间。必须借助于原子核自身的飞行运动学尺度进行超快时间测量。

\subsection{微观柱塞反冲距离多普勒频移法（RDDS / Plunger）}

用于测量寿命在 $1 \sim 500\text{ ps}$ 范围内的态。其核心装置是高精度的压电陶瓷柱塞（Plunger）：
\begin{enumerate}
    \item 将厚度约 $1\text{ mg/cm}^2$ 的超薄金属靶与一块重金属（金或铅）衬底平行安放，利用激光干涉系统将两者的空间物理间距 $d$ 精确控制在数微米至数毫米之间；
    \item 熔合反应产生的复合核以速度 $v = \beta c$ 冲出靶薄膜，飞向衬底；
    \item 若核在飞行途中退激，光子在全速度 $v$ 下发射，被探测器接收为完全频移的全多普勒峰（Shifted Peak, $I_s$）；若核在打入衬底完全停滞后才退激，光子在静止状态下发射，表现为未频移峰（Unshifted Peak, $I_u$）。
\end{enumerate}

两峰强度比值直接由简单的指数衰减定律给出：
\begin{equation}
    R(d) = \frac{I_u(d)}{I_s(d) + I_u(d)} = \exp\left( -\frac{d}{v \tau} \right)
    \label{eq:rdds_decay}
\end{equation}

在真实多能级级联衰变网络中，上述简单单指数衰减公式受到复杂的上能级直接馈入与侧翼未知馈入（Side Feeding）的严重污染。为了完全摆脱对未知侧翼馈入模型的唯象假定，A. Dewald 等人发展了著名的\term{微分衰减曲线法}（Differential Decay Curve Method, DDCM）：
\begin{equation}
    \tau_i(d) = \frac{I_u^{(i)}(d) - \sum_{j} b_{ji} I_u^{(j)}(d)}{v \cdot \frac{d}{dd} I_s^{(i)}(d)}
    \label{eq:ddcm_formula}
\end{equation}
其中 $b_{ji}$ 为已知母能级 $j$ 向待测能级 $i$ 跃迁的直接分支比。

式~(\ref{eq:ddcm_formula})表明：在任意间距 $d$ 处，分子为待测态静止峰强度减去所有已知上级联母态贡献，分母为频移峰面积对飞行距离的导数；求得的寿命 $\tau_i$ 在全部有效拟合距离范围内呈现为一个完美的水平平台。DDCM 方法在数学上彻底消除了对未观测高能侧翼馈入时间（Feeding Time）的任何模型依赖，成为当今国际反冲距离柱塞寿命测量的最高黄金准则。

\subsection{多普勒频移衰减法（DSAM）}

对于寿命更短的极高自旋态（$\tau \le 1\text{ ps}$），反冲核尚未飞出靶体或厚衬底，在减速慢化过程中就已发射光子。此时采用\term{多普勒频移衰减法}（Doppler Shift Attenuation Method, DSAM）：
\begin{itemize}
    \item 利用 Lindhard-Scharff-Schiott（LSS）理论与电子/核阻止本领精确模拟反冲核在固体介质中的速度衰减曲线 $v(t)$；
    \item 光子发射谱线产生连续的非对称展宽线形（Line-Shape）；
    \item 提取衰减因子 $F(\tau) = \frac{\langle \Delta E_D \rangle}{\Delta E_{D,\text{max}}} = \frac{1}{v_0 \tau} \int_0^\infty v(t) e^{-t/\tau} dt$，通过全波形线形拟合精确反演亚皮秒寿命。
\end{itemize}

\section{能级纲图构建原则与物理自洽性校验}
\label{sec:level_scheme_construction}

在束实验数据的终极科学成果是构建一份拓扑完整、逻辑自洽的\term{原子核能级纲图}（Level Scheme）。构建纲图决不能主观臆测，必须恪守以下四大约束定则：
\begin{enumerate}
    \item \textbf{Ritz 能量加和组合定则}：两条平行的并联跃迁路线，其跃迁能量求和必须在实验误差范围内严格守恒：$E_{\gamma 1} + E_{\gamma 2} = E_{\gamma 3} + E_{\gamma 4}$；
    \item \textbf{强度平衡原理（Intensity Balance）}：对于任何一个非同核异能态的激发能级，在稳定束流条件下，汇聚进入该能级的全部跃迁总强度（考虑内转换电子修正：$I_{\text{tot}} = I_\gamma (1 + \alpha)$），必须\textbf{严格大于等于}从该能级向下退激分支的总强度。如果发现出射强度大于入射强度，说明必然存在未被观测到的高自旋分支馈入（Side Feeding）；
    \item \textbf{符合门反向对称性}：如果门开在 $\gamma_1$ 上能清晰看到 $\gamma_2$，那么反向门开在 $\gamma_2$ 上必须百分之百确定看到 $\gamma_1$；任何单向单侧符合均视为假符合或被近邻杂质峰污染；
    \item \textbf{物理自洽性与单粒子极限约束}：跃迁几率绝不能超出 Weisskopf 单粒子估计上限。若某一条偶极跃迁反推算出的跃迁强度超过数十个单粒子 Weisskopf 单位，说明该跃迁的自旋指认或多极性指认存在颠覆性错误。
\end{enumerate}
